\begin{abstract}
Explanation heatmaps are often offered as evidence that a skin-lesion classifier attends to the lesion.
We ask whether two standard benchmarks for such heatmaps, deletion/insertion faithfulness and
localization against expert lesion masks, measure the model at all. On HAM10000 we fine-tune a CNN
(EfficientNet-B0) and a Vision Transformer (ViT-B/16) under one protocol with three training seeds and a
fixed lesion-grouped split, and score every model-explainer system on the same \todo{1,002} test images
and the same checkpoints. We pair each architecture's customary explainer (Grad-CAM, attention rollout)
with one shared, class-specific explainer (Integrated Gradients) and add two model-free controls: a
random pixel order and a fixed centered heatmap. The centered heatmap, which never looks at the image,
reaches an IoU of \todo{X} and an exact pointing-game rate of \todo{X}, \todo{[above/below]} Grad-CAM
(\todo{X}, \todo{X}) and attention rollout (\todo{X}, \todo{X}). \todo{One sentence on faithfulness vs.\
random and center-prior orderings.} \todo{One sentence on whether the CNN/ViT gap persists with the shared
explainer.} Because HAM10000 lesions are typically centered and cover about \todo{X}\% of the image,
localization scores on this dataset should be reported relative to a center prior, and comparisons
across architectures should hold the explainer fixed. Code will be made publicly available.
\end{abstract}
